% from ega1-8-fr.tex
\oldpage[I]{166}
\begin{proposition}[8.2.2]
\label{I.8.2.2-fr}
Soient $X$ un préschéma intègre, $R$ son corps des fonctions rationnelles.
Pour que $X$ soit un schéma, il faut et il suffit que la relation
«~$\mathscr{O}_x$ et $\mathscr{O}_y$ sont apparentés~»
(\hyperref[I.8.1.4-fr]{8.1.4}) entre points $x$, $y$ de $X$ implique $x=y$.
\end{proposition}

Supposons cette condition vérifiée, et montrons que $X$ est séparé. Soient
$U$ et $V$ deux ouverts affines distincts de $X$, $A$ et $B$ leurs anneaux,
identifiés à des sous-anneaux de $R$~; $U$ (resp. $V$) s'identifie donc
(\hyperref[I.8.1.2-fr]{8.1.2}) à $L(A)$ (resp. $L(B)$), et l'hypothèse
entraîne (\hyperref[I.8.1.5-fr]{8.1.5}) que si $C$ est le sous-anneau de $R$
engendré par $A\cup B$, $W=U\cap V$ s'identifie à
$L(A)\cap L(B)=L(C)$. En outre, on sait
(\cite{I-1}, p.~5-03, prop.~4 \emph{bis}) que tout sous-anneau $E$ de $R$ est
égal à l'intersection des anneaux locaux appartenant à $L(E)$~; $C$
s'identifie donc à l'intersection des anneaux $\mathscr{O}_z$ pour $z\in W$,
autrement dit (\hyperref[I.8.2.1.1-fr]{8.2.1.1}) à
$\Gamma(W,\mathscr{O}_X)$. Considérons alors le sous-préschéma induit par $X$
sur $W$~; à l'homomorphisme identique
$\varphi:C\to\Gamma(W,\mathscr{O}_X)$ correspond
(\hyperref[I.2.2.4-fr]{2.2.4}) un morphisme
$\Phi=(\psi,\theta):W\to\operatorname{Spec}(C)$~; nous allons voir que $\Phi$
est un \emph{isomorphisme} de préschémas, d'où résultera que $W$ est un ouvert
\emph{affine}. L'identification de $W$ à
$L(C)=\operatorname{Spec}(C)$ montre que $\psi$ est \emph{bijective}. D'autre
part, pour tout $x\in W$, $\theta_x^\sharp$ est l'injection
$C_\mathfrak r\to\mathscr{O}_x$, si
$\mathfrak r=\mathfrak m_x\cap C$, et par définition $C_\mathfrak r$ est
identifié à $\mathscr{O}_x$, donc $\theta_x^\sharp$ est bijective. Il reste
donc à voir que $\psi$ est un \emph{homéomorphisme}, autrement dit que pour
toute partie fermée $F\subset W$, $\psi(F)$ est fermé dans
$\operatorname{Spec}(C)$. Or, $F$ est la trace sur $W$ d'une partie fermée de
$U$, de la forme $V(\mathfrak a)$, où $\mathfrak a$ est un idéal de $A$~;
montrons que $\psi(F)=V(\mathfrak a C)$, ce qui prouvera notre assertion. En
effet, les idéaux premiers de $C$ contenant $\mathfrak a C$ sont les idéaux
premiers de $C$ contenant $\mathfrak a$, donc les idéaux de la forme
$\psi(x)=\mathfrak m_x\cap C$ où
$\mathfrak a\subset\mathfrak m_x$ et $x\in W$~; comme
$\mathfrak a\subset\mathfrak m_x$ équivaut à
$x\in V(\mathfrak a)=W\cap F$ pour $x\in U$, on a bien
$\psi(F)=V(\mathfrak a C)$.

Il s'ensuit que $X$ est séparé, car $U\cap V$ est affine et son anneau $C$ est
engendré par la réunion $A\cup B$ des anneaux de $U$ et $V$
(\hyperref[I.5.5.6-fr]{5.5.6}).

Inversement, supposons $X$ séparé, et soient $x$, $y$ deux points de $X$ tels
que $\mathscr{O}_x$ et $\mathscr{O}_y$ soient apparentés. Soit $U$ (resp.
$V$) un ouvert affine contenant $x$ (resp. $y$), d'anneau $A$ (resp. $B$)~;
on sait alors que $U\cap V$ est affine et que son anneau $C$ est engendré par
$A\cup B$ (\hyperref[I.5.5.6-fr]{5.5.6}). Si
$\mathfrak p=\mathfrak m_x\cap A$, $\mathfrak q=\mathfrak m_y\cap B$, on a
$A_\mathfrak p=\mathscr{O}_x$, $B_\mathfrak q=\mathscr{O}_y$, et comme
$A_\mathfrak p$ et $B_\mathfrak q$ sont apparentés, il existe un idéal premier
$\mathfrak r$ de $C$ tel que $\mathfrak p=\mathfrak r\cap A$,
$\mathfrak q=\mathfrak r\cap B$ (\hyperref[I.8.1.4-fr]{8.1.4}). Mais alors il
existe un point $z\in U\cap V$ tel que
$\mathfrak r=\mathfrak m_z\cap C$ puisque $U\cap V$ est affine, et on a
évidemment $x=z$, et $y=z$, d'où $x=y$.

\begin{corollary}[8.2.3]
\label{I.8.2.3-fr}
Soient $X$ un schéma intègre, $x$, $y$ deux points de $X$. Pour que
$x\in\overline{\{y\}}$, il faut et il suffit que
$\mathscr{O}_x\subset\mathscr{O}_y$, autrement dit que toute fonction
rationnelle définie en $x$ soit définie en $y$.
\end{corollary}

La condition est évidemment nécessaire puisque le domaine de définition
$\delta(f)$ d'une fonction rationnelle $f\in R$ est ouvert~; montrons qu'elle
est suffisante. Si $\mathscr{O}_x\subset\mathscr{O}_y$, il existe un idéal
premier $\mathfrak p$ de $\mathscr{O}_x$ tel que $\mathscr{O}_y$ domine
$(\mathscr{O}_x)_\mathfrak p$ (\hyperref[I.8.1.3-fr]{8.1.3})~; or
(\hyperref[I.2.4.2-fr]{2.4.2}) il existe $z\in X$ tel que
$x\in\overline{\{z\}}$ et que
$\mathscr{O}_z=(\mathscr{O}_x)_\mathfrak p$~; comme $\mathscr{O}_z$ et
$\mathscr{O}_y$ sont apparentés, on a $z=y$ par
(\hyperref[I.8.2.2-fr]{8.2.2}), d'où le corollaire.

\begin{corollary}[8.2.4]
\label{I.8.2.4-fr}
Si $X$ est un schéma intègre, l'application $x\mapsto\mathscr{O}_x$ est
injective~; autrement dit, si $x$, $y$ sont deux points distincts de $X$, il
existe une fonction rationnelle définie en l'un de ces points et non en
l'autre.
\end{corollary}

